\documentclass[12pt,a4paper]{article}
% ============================================================
%  Structural Persistence Series — Shared Preamble
% ============================================================
\usepackage{fontspec}
\setmainfont{Hiragino Mincho ProN}
\setsansfont{Hiragino Sans}
\setmonofont{Menlo}

\usepackage{iftex}
\ifXeTeX
  \XeTeXlinebreaklocale "ja"
  \XeTeXlinebreakskip=0pt plus 1pt minus 0.1pt
\fi

\usepackage[margin=22mm]{geometry}
\usepackage{amsmath,amssymb,amsthm}
\usepackage{xcolor}
\usepackage{graphicx}
\usepackage{hyperref}
\usepackage{booktabs}
\usepackage{fancyhdr}
% enumitem not available in this TeX installation

% --- Colours ------------------------------------------------
\definecolor{PaperAccent}{HTML}{1F4E79}
\definecolor{PaperGray}{HTML}{51606F}

\hypersetup{
  colorlinks=true,
  linkcolor=PaperAccent,
  urlcolor=PaperAccent,
  citecolor=PaperAccent,
  pdflang=ja
}
\urlstyle{same}

% --- Typography ---------------------------------------------
\setlength{\parindent}{1.5em}
\setlength{\parskip}{0pt}
\setlength{\emergencystretch}{3em}
\linespread{1.08}
\setlength{\abovedisplayskip}{8pt plus 2pt minus 4pt}
\setlength{\belowdisplayskip}{8pt plus 2pt minus 4pt}
\setlength{\abovedisplayshortskip}{6pt plus 2pt minus 2pt}
\setlength{\belowdisplayshortskip}{6pt plus 2pt minus 2pt}

% --- Section label names ------------------------------------
\renewcommand{\abstractname}{要旨}

% --- Theorem-like environments ------------------------------
\theoremstyle{plain}
\newtheorem{proposition}{命題}[section]
\newtheorem{lemma}[proposition]{補題}
\newtheorem{corollary}[proposition]{系}

\theoremstyle{definition}
\newtheorem{definition}{定義}[section]
\newtheorem{assumption}{条件}[section]

\theoremstyle{remark}
\newtheorem*{remark}{注}

% --- Running header -----------------------------------------
\pagestyle{fancy}
\fancyhf{}
\renewcommand{\headrulewidth}{0.4pt}
\fancyhead[L]{\small\textcolor{PaperGray}{\nouppercase{\leftmark}}}
\fancyhead[R]{\small\textcolor{PaperGray}{\thepage}}
\fancypagestyle{plain}{%
  \fancyhf{}
  \renewcommand{\headrulewidth}{0pt}
  \fancyfoot[C]{\small\textcolor{PaperGray}{\thepage}}
}

% --- Title block macro --------------------------------------
\newcommand{\PaperTitleBlock}[3]{%
  % #1 = title, #2 = subtitle, #3 = date
  \thispagestyle{plain}%
  \begin{center}
  {\LARGE\bfseries #1\par}
  \vspace{0.4em}
  {\normalsize #2\par}
  \vspace{1.0em}
  {\large Akihito Sunagawa\par}
  \vspace{0.3em}
  {\normalsize #3\par}
  \end{center}
  \vspace{1.6em}
}

% Legacy 2-arg version (backward compat)
\newcommand{\PaperTitleBlockSimple}[2]{%
  \PaperTitleBlock{#1}{}{#2}
}


\begin{document}

\PaperTitleBlock
  {構造持続の最小形式}
  {--- 制約の蓄積と構造損失 ---}
  {2026年4月}

\begin{abstract}
構造が失われるのは、資源が尽きるときだけではない。資源がなお残っていても、制約の蓄積によって構造を維持できる状態の領域が縮小し、ついにはその構造を保てる状態が残らなくなることがある。本稿は、この事実を表す最小形式を与える。系に対し、その構造を維持できる状態の集合を置き、制約の蓄積をその集合の逐次縮小として表す。各段階の損失を対数比で定義すると、損失は加法的に累積し、残存可能性は指数型で表される。したがって、指数型は追加仮定ではなく、この表現の帰結である。さらに、有効維持資源を導入し、構造持続ポテンシャル
\begin{equation*}
S = Me^{-L}
\end{equation*}
を定義する。この形式は、資源不足と構造的維持可能性の減少とを切り分けて記述する。
\end{abstract}

% =============================================================
\section{はじめに}
% =============================================================

本稿の問いは、構造は何によって持続し、何によって失われるのか、である。
構造の喪失を資源の減少だけで捉えるなら、この問いは十分には答えられない。
むしろ問うべきなのは、構造を維持できる状態の領域が何によって縮小するのか、である。

本稿の目的は、この事実を表す最小形式を与えることである。以下では、系に対して構造維持可能な状態集合を置き、制約の蓄積をその逐次縮小として表す。各段階の損失を対数比で定義することにより、損失の加法性と残存可能性の指数表現が導かれることを示す。さらに、有効維持資源を導入することで、資源の不足と構造を維持できる状態の減少という二つの側面を切り分けて記述する。

% =============================================================
\section{設定}
% =============================================================

系 $X$ を固定する。本稿でいう「構造」とは、系において維持すべき関係・機能・同一性の総称であり、その具体的内容は個別領域に委ねる。初期時点において、その構造を維持できる状態の集合を $V^{(0)}$ とする。制約が順に蓄積するとき、構造を維持できる状態の領域は各段階で縮小し、
\begin{equation}\label{eq:chain}
V^{(0)} \supseteq V^{(1)} \supseteq \cdots \supseteq V^{(n)}
\end{equation}
を得るものとする。本稿は、制約追加が支配的であり、学習や修復による再拡大を無視できる局面を対象とする。

各 $V^{(i)}$ は、ある可測空間上の可測集合であり、その上に固定された有限測度 $m$ を入れる。$m(V^{(i)})$ は、その段階でなお残っている構造維持可能性の大きさを表す。$m$ は個数でも体積でも確率重みでもよいが、本稿では、同一の系の同一表現の上で段階を通じて比較可能な一つの測度であることのみを要請する。なお、$m(V^{(i)})$ は系が現に占めている状態ではなく、その段階でなお構造維持と両立しうる状態の余地を表す。

% =============================================================
\section{各段階の損失}
% =============================================================

各段階 $i$ における損失 $l_i$ を
\begin{equation}\label{eq:stage-loss}
l_i = -\ln\!\left(\frac{m(V^{(i)})}{m(V^{(i-1)})}\right)
\end{equation}
によって定義する。ただし $m(V^{(i-1)}) > 0$ とする。もし $m(V^{(i)}) = 0$ ならば $l_i = +\infty$ とおく。

$l_i$ は、負荷の大きさそのものではなく、その段階で構造を維持できる状態の集合がどれだけ縮小したかを表す量である。したがって、これは、構造を維持しうる可能性の量が段階ごとに削られていくことを表す対数的損失であり、その意味で、構造維持可能性に関する自由度の喪失として読むことができる。

累積損失 $L_n$ を
\begin{equation}\label{eq:cumulative-loss}
L_n = \sum_{i=1}^{n} l_i
\end{equation}
と定める。

% =============================================================
\section{指数表現}
% =============================================================

\begin{proposition}\label{prop:exp}
$V^{(0)} \supseteq V^{(1)} \supseteq \cdots \supseteq V^{(n)}$ を有限測度 $m$ の下での可測集合列とし、
\begin{equation*}
l_i = -\ln\!\left(\frac{m(V^{(i)})}{m(V^{(i-1)})}\right),
\qquad
L_n = \sum_{i=1}^{n} l_i
\end{equation*}
と定める。このとき
\begin{equation}\label{eq:exp-form}
m(V^{(n)}) = m(V^{(0)})\, e^{-L_n}
\end{equation}
が成り立つ。
\end{proposition}

\begin{proof}
定義より
$e^{-l_i} = m(V^{(i)}) / m(V^{(i-1)})$
であるから、
\begin{align*}
e^{-L_n}
&= \prod_{i=1}^{n} e^{-l_i}
= \prod_{i=1}^{n} \frac{m(V^{(i)})}{m(V^{(i-1)})}
= \frac{m(V^{(n)})}{m(V^{(0)})}.
\end{align*}
よって $m(V^{(n)}) = m(V^{(0)})\, e^{-L_n}$ が従う。
\end{proof}

したがって指数型は追加仮定ではなく、状態集合の乗法的縮小を対数比で表したことの直接の帰結である。その意味で指数型それ自体は経験法則ではなく、経験的内容は $V^{(i)}$ と $m$ の具体化、および各制約に対応する縮小率の差異に属する。本形式が単なる記述以上の予測力を持つのは、制約ごとの縮小率 $l_i$ が異なりうることを通じて、制約の質が量に優先しうるという順序予測を生成するときである。この予測が個別ドメインで検証可能かどうかが、本形式の経験的試練となる。

本稿の主張がどこにあるかを明示する。本形式において、対数比で定義された段階損失の加法的累積から指数型が現れること自体は、既知の数学的構造に属する。望遠鏡積による指数表現は、生存解析におけるハザード関数の累積、統計力学における分配関数の構造、情報理論における容量境界と形式的に同型であり、指数型そのものに数学的新規性はない。

本稿の賭けは、この既知の骨格が、知能的・推論的・学習的な構造の崩壊に対して、有効な概念写像を与えうるかどうか、にある。したがって、本形式の経験的内容は指数型そのものにはなく、各ドメインにおいて何を構造維持可能状態とみなすか、どの測度で比較するか、どの制約がどれだけの縮小率を持つか、という具体化の側に属する。

この点を認めることは、本形式を弱めるのではなく、むしろ勝負所を明確にする。すなわち、本形式が試されるべきなのは、「指数型が成り立つか」ではなく、「この写像手順が、異なるドメインにおいて、基準モデルに対する追加の予測力を生むか」である。

% =============================================================
\section{構造持続ポテンシャル}
% =============================================================

命題~\ref{prop:exp}は、構造を維持できる状態集合の残存量がどのように表されるかを与える。しかし、構造の持続は、状態の広がりだけで決まるとは限らない。実際の系では、その構造を保つために利用できる資源もまた重要である。そこで、資源要因と構造維持可能性の縮小要因とを独立な寄与として分離して表す最小の定式化として、命題~\ref{prop:exp}の幾何学的縮小とは別に、有効維持資源を $M > 0$ として導入し、構造持続ポテンシャル $S$ を
\begin{equation}\label{eq:S}
S = Me^{-L_n}
\end{equation}
と定める。これは命題~\ref{prop:exp}の直接の帰結ではなく、資源要因を組み込んだ拡張的定義である。

$M$ は $m(V^{(0)})$ と同一視される必要はない。$m(V^{(0)})$ は構造維持可能状態の初期的広がりであるのに対し、$M$ は系がその構造を維持するために利用できる有効資源の大きさを表す。個別の系において、それはエネルギー、時間、冗長性、回復可能性、計算余地、あるいは流れを維持する能力として具体化されうる。本稿では、その内訳には立ち入らない。

最も単純には、構造の持続はある閾値 $S_c$ に対して
\begin{equation}\label{eq:threshold}
S < S_c
\end{equation}
となるとき失われる。このとき観測される現象は、崩壊、機能不全、相転移、自己同一性の喪失など、系に応じて異なる形をとる。

% =============================================================
\section{この形式が含意すること}
% =============================================================

\subsection{制約の質は量に優先しうる}

各段階の寄与は、制約の数ではなく、状態集合の縮小率によって測られる。したがって、少数の制約でも、構造を維持できる状態を大きく減少させるならば支配的になりうる。逆に、多数の制約でも、その縮小率が小さければ、累積損失は大きくならない。したがって本形式は、制約の個数よりも、どの制約がより大きく維持可能状態集合を削るかに関する順序予測を与える。

\subsection{資源と損失は異なる側面をなす}

$S = M\exp(-L_n)$ において、$M$ と $L_n$ は構造持続の異なる側面を表す。一方は維持資源の大きさであり、他方は構造を維持できる状態の縮小である。いずれか一方が十分に悪化すれば、構造持続ポテンシャルは失われる。したがって、資源が残っていても構造を維持できる状態が失われる場合があり、逆に構造を維持できる状態がなお残っていても、資源不足によって維持が失われる場合がある。

\subsection{臨界領域が現れうる}

$S$ が閾値 $S_c$ に近い領域では、小さな追加損失でも構造の持続が失われうる。したがって、閾値的な挙動を示す系では、緩やかな劣化ではなく、急峻な転移に近い挙動が現れうる。

% =============================================================
\section{例}
% =============================================================

この形式は、個別の系に応じて異なる対象に対応づけられる。

\textbf{推論系}では、$V^{(i)}$ は、整合性を保ったまま推論を完了できる状態または経路の集合に対応する。不整合や未整理の上書きは、その集合を縮小させる制約として現れる。

\textbf{SAT} では、$V^{(i)}$ は、その段階でなお残っている充足割当ての集合に対応する。各制約はこの集合を縮小させ、累積損失は解空間の収縮として読まれる。個別領域において制約型ごとの損失が評価可能な場合には、本形式は比率予測へ接続しうる。SAT はその具体例である。詳細は別稿に譲る。

\textbf{散逸構造や生物系}では、$V^{(i)}$ は、境界維持、流れ、代謝、温度勾配、化学ポテンシャル差などの条件と両立する状態の集合に対応する。このとき構造持続は、静止した秩序としてではなく、流れの中で維持される自己同一性として現れる。

本稿の主張は、これらを同一視することにはない。主張は、それらがいずれも「構造を維持できる状態集合の縮小」という同じ最小形式のもとで記述可能である、という一点にある。

% =============================================================
\section{限界}
% =============================================================

本稿は、すべての系に対して一つの測度 $m$、一つの資源量 $M$、一つの閾値 $S_c$ を与えるものではない。また、個別の系において $L_n$ が一意に測定できることも主張しない。さらに、制約の追加が常にあらゆる表現において縮小方向にのみ働くことを、一般定理として主張するものでもない。

とくに、本形式は制約蓄積が支配的な局面を記述するものであり、学習・修復・冗長化など、構造を維持できる状態集合を再拡大する過程は明示的に含まれていない。したがって、本形式は、散逸構造を含む広い範囲の系に対して持続条件の骨格を与えうるが、流れによる自己維持や再編成のような動的過程そのものを記述するものではない。そのような動的過程を含む理論は、今後の拡張として別途扱う必要がある。

測度 $m$ の具体化は個別領域に依存する。SAT では充足割当ての数え上げとして自然に定まるが、推論系のように直接の測定が困難な領域では、観測指標を介した間接的な接近が必要になる。このとき本形式は、指標の探索と検証のための枠組みとして機能する。すなわち、本形式が与える順序予測や方向予測に合致する指標は妥当な代理指標として残り、合致しない指標は棄却される。したがって、測度の直接測定が困難であることは、本形式の適用を不可能にするのではなく、代理指標の発見を構造化された経験的課題として位置づける。

本稿が与えるのは、構造持続を記述するための条件つきの最小形式である。すなわち、ある系について構造を維持できる状態集合が定まり、制約の蓄積をその集合の縮小として表せるならば、損失は対数比として定義でき、残存可能性は指数型で表される、ということである。

% =============================================================
\section{結論}
% =============================================================

構造の持続は、単なる資源残量の問題ではない。構造維持条件と両立する状態がなお存在するかどうかの問題でもある。本稿では、その事実を表す最小形式を与えた。構造を維持できる状態集合の逐次縮小を対数比で表すと、損失は加法的に累積し、残存可能性は指数型で表される。指数型は仮定ではなく、この表現の帰結である。さらに、有効維持資源を導入することで、構造持続ポテンシャル
\begin{equation*}
S = Me^{-L}
\end{equation*}
を得た。この形式は、資源の不足と、構造を維持できる状態の減少とを切り分けつつ、構造持続の喪失を一つの形で記述する。

本稿が与えるのは、個別理論に先立つ最小の骨格である。測度の選択、資源量の定義、閾値の同定、実証、および応用は別稿に譲る。

\section*{付記}

本文では簡潔さのために $L_n$ と書いたが、文脈が明らかなときは単に $L$ と書いてよい。したがって、中心式は $S = Me^{-L}$ である。

% =============================================================
\begin{thebibliography}{9}
\bibitem{paper2} Sunagawa, A. 構造持続の条件つき導出. 2026.
\bibitem{paper3} Sunagawa, A. 論理一貫性維持の構造持続的記述. 2026.
\bibitem{paper4} Sunagawa, A. 前提更新を伴う継続学習における構造的忘却. 2026.
\bibitem{paper5} Sunagawa, A. 持続知能の混成構成. 2026.
\end{thebibliography}

\end{document}
